\scanpage{46}
\mspage{- 23 -}
\margin{n° 380}

§ 2. \uuline{Calcul fonctionnel holomorphe}.

1. \uline{Homomorphisme} $\mathcal O(\operatorname{Sp}(x))\longrightarrow A$

\uline{Les deux lemmes suivants figureront au fascicule de Résultats ou autre part}.

LEMME 1.- \uline{Soient} $K$ \uline{un compact de} $\mathbb C$ \uline{et} $V$ \uline{un voisinage ouvert de} $K$. \uline{Il existe une mesure} $d\rho$ \uline{à support compact contenu dans} $V\cap(\mathbb C-K)$ \uline{et un voisinage ouvert} $W$ \uline{de} $K$ \uline{contenu dans} $V$ \uline{tels que}, \uline{pour toute fonction holomorphe} $f$ \uline{sur} $V$ \uline{à valeurs dans un espace de Banach} $E$ \uline{et tout} $z\in W$, \uline{on ait}
\[
f(z)=\int\frac{f(t)}{z-t}\,d\rho(t).
\]

LEMME 2.- \uline{Soient} $L$ \uline{un compact de} $\mathbb C$ \uline{et} $U$ \uline{un voisinage ouvert de} $L$. \uline{Il existe une mesure} $d\rho$ \uline{à support compact contenu dans} $U\cap(\mathbb C-L)$ \uline{et un voisinage ouvert} $W$ \uline{de} $\mathbb C-U$ \uline{tels que}, \uline{pour toute fonction holomorphe} $F$ \uline{sur} $\mathbb C-K$, \uline{nulle à l'infini}, \uline{à valeurs dans un espace de Banach} $E$ \uline{et tout} $z\in W$
\[
F(z)=\int\frac{F(t)}{z-t}\,d\rho(t).
\]

LEMME 3.- \uline{Soit} $K$ \uline{un compact dans} $\mathbb C$ \uline{et soit} $d\unsure{\mathcal D|\nu}$ \uline{une mesure à support compact contenu dans} $\Omega=\mathbb C-K$ \uline{telle que} $\int\frac{d\unsure{\mathcal D|\nu}(t)}{z-t}=0$ \uline{pour tous les points} $z$ \uline{d'un voisinage de} $K$. \uline{Si} $F$ \uline{est une fonction holomorphe sur} $\Omega$ \uline{à valeurs dans un espace de Banach} $E$ \uline{et si} $F$ \uline{est nulle à l'infini}, \uline{alors} $\int F(t)\,d\unsure{\mathcal D|\nu}(t)=0$.

Soit en effet $U$ un voisinage ouvert relativement compact de $K$ ne rencontrant pas le support de $d\unsure{\mathcal D|\nu}$ et tel que $\int\frac{d\unsure{\mathcal D|\nu}(t)}{z-t}=0$ pour tout $z\in U$. Puisque $F$ est nulle à l'infini, il existe une mesure $d\rho$ à support compact contenu dans $U\cap\Omega$ telle que pour tout $z\in\mathbb C-U$ on ait
\[
F(z)=\int\frac{F(t)\,d\rho(t)}{z-t}.
\]
(Lemme 2). On a donc, d'après Fubini :

\margin{56}
